\section{Conclusiones}

\begin{enumerate}

\item \textbf{DistilBERT es la elección por defecto para estas tareas.}
Conserva entre el \SI{98,9}{\percent} y el \SI{99,5}{\percent} de la
\textit{accuracy} de BERT con el \SI{61}{\percent} de los parámetros, el
\SI{62}{\percent} de la memoria y la mitad de la latencia. La pregunta
relevante no es si BERT es mejor --- lo es, por poco --- sino si ese margen
justifica duplicar el costo de inferencia. En AG News no lo justifica en
absoluto; en SST-2, donde la brecha llega a \num{1,95} puntos con la cabeza
lineal, podría hacerlo si la aplicación es sensible a esos puntos.

\item \textbf{La brecha entre los modelos depende de la tarea, no es una
constante.} Va de \num{0,11} puntos en AG News a \num{1,95} en SST-2. La
clasificación de temas se resuelve en buena medida con vocabulario; el
análisis de sentimiento sobre frases cortas, con negación y contraste, exige
la composición que aportan las capas adicionales. Extrapolar un único número
de ``retención de desempeño'' entre tareas sería un error.

\item \textbf{La capacidad que importa está en el \textit{encoder}, no en la
cabeza.} Cinco arquitecturas de cabeza --- de un clasificador lineal a un MLP
de dos capas ocultas --- caben en \num{0,0014} de F1, mientras que congelar el
\textit{transformer} cuesta entre \num{3,4} y \num{8,4} puntos. El corolario
práctico es que el esfuerzo de diseño debe ir a qué se adapta y con qué datos,
no a la arquitectura del clasificador.

\item \textbf{La medición de eficiencia exige condiciones controladas.} La
latencia registrada al final de cada entrenamiento sugería que BERT era más
rápido que DistilBERT en SST-2. Medidos en la misma rejilla, mismo proceso y
misma GPU, DistilBERT resulta \num{2}$\times$ más rápido en los nueve puntos.
La diferencia entre las dos mediciones no está en los modelos sino en el
protocolo, y es un recordatorio de que una métrica de sistema sin condiciones
explícitas no es un resultado.

\item \textbf{Congelar el \textit{encoder} es un intercambio, no un error.}
Entrena \num{0,6}~M de parámetros en lugar de \num{66,4}~M y tarda menos de la
mitad. En AG News el costo es de \num{3,4} puntos de F1, que puede ser
aceptable si el presupuesto de entrenamiento es el factor limitante.

\end{enumerate}
